\documentclass[12pt]{article}
\usepackage[utf8]{inputenc}
\usepackage[T1]{fontenc}
\usepackage[margin=1in]{geometry}
\usepackage{setspace}
\usepackage{xcolor}
\usepackage[colorlinks=true, urlcolor=blue!60!black, linkcolor=blue!60!black]{hyperref}
\setlength{\parindent}{2em}
\setlength{\parskip}{1.25em}
\renewcommand{\baselinestretch}{1.0}
\begin{document}
\begin{center}
{\LARGE \textit{National Pork Producers Council v. Ross}, 598 U.S. 356 (2023)\par}
\vspace{0.8em}
{\large Case Brief\par}
\end{center}
\vspace{1em}

\section{Memory Jogger}
Affirming dismissal of a Dormant Commerce Clause challenge to California's Proposition 12, the Supreme Court unanimously rejected an ``almost per se'' rule against state laws with extraterritorial effects and fractured 5--4 in upholding California's nondiscriminatory ban on in-state sales of pork from cruelly confined breeding pigs---leaving uniform national rules to Congress under its affirmative Commerce Clause and Supremacy Clause preemption powers.

\section{Detailed Case Facts}
In November 2018, roughly 63\% of participating California voters approved Proposition 12, a ballot initiative amending the California Health and Safety Code (\S\S~25990--25994) to revise in-state sales standards for eggs and establish new standards for veal and whole pork meat (\href{https://www.courtlistener.com/opinion/10049695/national-pork-producers-council-v-ross/}{CourtListener}; \href{https://www.law.cornell.edu/supremecourt/text/21-468}{Cornell LII}). Specifically, Proposition 12 prohibited the in-state commercial sale of whole pork meat derived from a breeding pig (or its immediate offspring) ``confined in a cruel manner,'' defined as preventing the animal from ``lying down, standing up, fully extending [its] limbs, or turning around freely,'' along with minimum square-footage standards (Cal. Health \& Safety Code \S\S~25990(b)(2), 25991(e)(1)) (\href{https://www.law.cornell.edu/supremecourt/text/21-468}{Cornell LII}; \href{https://www.gwlr.org/national-pork-producers-council-v-ross-a-win-for-farmed-pigs/}{GW Law Review}). Proponents argued that Proposition 12 advanced animal welfare by ending 16-week gestation-crate confinement and protected consumer health by reducing foodborne-illness risks; opponents contended that existing confinement practices prevented pig-on-pig aggression, safeguarded herd health, and avoided higher consumer prices (\href{https://www.law.cornell.edu/supremecourt/text/21-468}{Cornell LII}).

Two national agricultural trade associations---the National Pork Producers Council (NPPC) and the American Farm Bureau Federation---sued Karen Ross, Secretary of the California Department of Food and Agriculture, and other state officials on behalf of their members (\href{https://www.law.cornell.edu/supremecourt/text/21-468}{Cornell LII}). The Humane Society of the United States and allied animal-welfare groups intervened as defendants (\href{https://www.law.cornell.edu/supct/cert/21-468}{Cornell LII Bulletin}). In their complaint, petitioners alleged that the U.S. pork industry is a \$20 billion-plus vertically integrated market in which different cuts from a single hog are shipped to different States, making segregation and tracing costly. Because California accounts for roughly 13\% of national pork consumption while importing nearly all of its pork from major producing States such as Iowa, Minnesota, Illinois, Indiana, and North Carolina, petitioners alleged that Proposition 12 would require \$290 million to \$348 million in capital expenditures, raise farm-level production costs by 9.2\% (over \$13 per pig), and practically force out-of-state producers nationwide to restructure their operations (\href{https://www.law.cornell.edu/supremecourt/text/21-468}{Cornell LII}). Petitioners conceded, however, that 28\% of the industry had already adopted some form of group housing, that Proposition 12 applied identically to in-state and out-of-state pork producers without purposeful discrimination, and that Congress---despite repeated lobbying by pork producers---had declined to enact federal legislation preempting state pork-confinement laws under the Commerce and Supremacy Clauses (Art. I, \S~8, cl. 3; Art. VI, cl. 2) (\href{https://www.law.cornell.edu/supremecourt/text/21-468}{Cornell LII}).

\section{Procedural History}
Petitioners filed suit on December 5, 2019, in the U.S. District Court for the Southern District of California, seeking declaratory and injunctive relief under the Dormant Commerce Clause on two theories: (1) an ``almost per se'' extraterritoriality doctrine, and (2) the balancing test of \textit{Pike v. Bruce Church, Inc.}, 397 U.S. 137 (1970) (\href{https://www.law.cornell.edu/supct/cert/21-468}{Cornell LII Bulletin}). The District Court dismissed the complaint under Federal Rule of Civil Procedure 12(b)(6) for failure to state a claim, 456 F. Supp. 3d 1201 (S.D. Cal. 2020) (\href{https://www.law.cornell.edu/supremecourt/text/21-468}{Cornell LII}).

On July 28, 2021, a unanimous panel of the U.S. Court of Appeals for the Ninth Circuit, in an opinion by Judge Ikuta, affirmed the dismissal, 6 F.4th 1021 (9th Cir. 2021), holding that upstream practical effects and increased compliance costs, without more, neither violated the narrow extraterritoriality doctrine nor established a substantial burden on interstate commerce under \textit{Pike} (\href{https://www.law.cornell.edu/supremecourt/text/21-468}{Cornell LII}). The Supreme Court granted certiorari on March 28, 2022 (Docket No. 21--468), heard oral argument on October 11, 2022 (with the United States participating as \textit{amicus curiae} in support of petitioners), and affirmed on May 11, 2023 (\href{https://supreme.justia.com/cases/federal/us/598/21-468/}{Justia}; \href{https://www.scotusblog.com/cases/national-pork-producers-council-v-ross/}{SCOTUSblog}).

\section{Judicial Votes}
The decision produced a famously cross-cutting ideological alignment and a classic Supreme Court voting paradox (\href{https://www.law.cornell.edu/supremecourt/text/21-468}{Cornell LII}; \href{https://www.scotusblog.com/cases/national-pork-producers-council-v-ross/}{SCOTUSblog}):

\begin{itemize}
\item \textbf{Judgment (5--4 to Affirm):} Justice Gorsuch announced the judgment of the Court affirming dismissal, joined in the judgment by Justices Thomas, Sotomayor, Kagan, and Barrett. Chief Justice Roberts and Justices Alito, Kavanaugh, and Jackson dissented from the judgment and would have vacated and remanded for \textit{Pike} balancing.
\item \textbf{Opinion of the Court---Parts I, II, III, IV--A, and V (5 Justices):} Justice Gorsuch, joined by Justices Thomas, Sotomayor, Kagan, and Barrett, rejected petitioners' ``almost per se'' extraterritoriality theory and reaffirmed that the antidiscrimination principle lies at the ``very core'' of Dormant Commerce Clause doctrine. On the rejection of the ``almost per se'' extraterritoriality rule, the Court was effectively \textbf{9--0}, as Chief Justice Roberts's partial dissent (joined by Alito, Kavanaugh, and Jackson, JJ.) expressly agreed with that conclusion.
\item \textbf{Part IV--B and Part IV--D---Incommensurability Bar to \textit{Pike} Balancing (3-Justice Plurality):} Justice Gorsuch, joined only by Justices Thomas and Barrett, concluded that courts cannot weigh incommensurable economic costs against noneconomic moral and health benefits. \textit{Crucially, six Justices rejected this view}: Justices Sotomayor and Kagan (concurring in part) plus Chief Justice Roberts and Justices Alito, Kavanaugh, and Jackson (concurring in part and dissenting in part) affirmed that \textit{Pike} balancing remains viable and judicially administrable.
\item \textbf{Part IV--C---No Substantial Burden Pleaded under \textit{Exxon}/\textit{Pike} (4-Justice Plurality):} Justice Gorsuch, joined by Justices Thomas, Sotomayor, and Kagan, concluded that under \textit{Exxon Corp. v. Governor of Maryland}, 437 U.S. 117 (1978), petitioners failed to plausibly plead a substantial burden on interstate commerce. \textit{Crucially, five Justices disagreed}: Justice Barrett (concurring in part) plus Chief Justice Roberts and Justices Alito, Kavanaugh, and Jackson concluded that petitioners \textit{did} plausibly allege a substantial burden.
\item \textbf{Separate Opinions:} Justice Sotomayor concurred in part (joined by Kagan, J.); Justice Barrett concurred in part (solo); Chief Justice Roberts concurred in part and dissented in part (joined by Alito, Kavanaugh, and Jackson, JJ.); and Justice Kavanaugh filed a separate opinion concurring in part and dissenting in part.
\end{itemize}

\section{Holding}
\subsection{Judgment}
The judgment of the Ninth Circuit affirming the Rule 12(b)(6) dismissal of petitioners' Dormant Commerce Clause complaint was \textbf{affirmed} (5--4) (\href{https://www.law.cornell.edu/supremecourt/text/21-468}{Cornell LII}).

\subsection{Rule Announced}
\begin{enumerate}
\item \textbf{Core Antidiscrimination Principle:} The antidiscrimination principle---forbidding state laws driven by economic protectionism designed to benefit in-state economic interests by burdening out-of-state competitors---lies at the ``very core'' of the Dormant Commerce Clause.
\item \textbf{No ``Almost Per Se'' Extraterritoriality Rule:} The Dormant Commerce Clause does not contain an ``almost per se'' rule invalidating nondiscriminatory state laws merely because they have the practical effect of influencing or altering out-of-state commercial behavior. Earlier cases such as \textit{Baldwin v. G.A.F. Seelig, Inc.}, 294 U.S. 511 (1935), \textit{Brown-Forman Distillers Corp. v. New York State Liquor Authority}, 476 U.S. 573 (1986), and \textit{Healy v. Beer Institute}, 491 U.S. 324 (1989), turned on purposeful economic protectionism in price-control and price-affirmation schemes.
\item \textbf{Narrowed \textit{Pike} Scope and Congressional Primacy:} While \textit{Pike} balancing survives and the courtroom door remains open to nondiscriminatory challenges, heartland \textit{Pike} cases serve primarily to smoke out hidden protectionism or protect instrumentalities of interstate transportation. Absent federal preemption enacted by Congress under Art. I, \S~8, cl. 3 and Art. VI, cl. 2, courts must exercise ``extreme caution'' before invalidating nondiscriminatory state regulations of in-state consumer-goods sales.
\end{enumerate}

\subsection{Holding vs. Dicta}
\begin{itemize}
\item \textbf{Binding Holdings:} Parts I, II, III, IV--A, and V of Justice Gorsuch's opinion command a five-Justice majority and bind all lower courts: there is no freestanding ``almost per se'' extraterritoriality bar, and \textit{Pike} operates primarily in tandem with the antidiscrimination principle.
\item \textbf{Fractured Plurality and \textit{Marks} Ambiguity on Part IV:} Because neither Part IV--B (3 Justices: incommensurability) nor Part IV--C (4 Justices: failure to plead a substantial burden under \textit{Exxon}) commanded a majority---and both rationales were rejected by majorities of the Court (6--3 and 5--4, respectively)---identifying the controlling \textit{Pike} rationale under \textit{Marks v. United States} is contested. Justice Kavanaugh contended in his separate opinion that Part IV--C is controlling precedent as the narrower ground, whereas Justice Gorsuch asserted in footnote 4 that the judgment depended equally on Parts IV--B and IV--C (\href{https://www.law.cornell.edu/supremecourt/text/21-468}{Cornell LII}).
\item \textbf{Dicta:} Justice Kavanaugh's discussion of potential future challenges under the Import-Export Clause, Privileges and Immunities Clause, and Full Faith and Credit Clause (issues neither pleaded nor presented by petitioners) is pure dicta, as is Justice Gorsuch's passing reference in Part II to originalist critiques questioning the Dormant Commerce Clause's textual home (\href{https://www.gwlr.org/national-pork-producers-council-v-ross-a-win-for-farmed-pigs/}{GW Law Review}; \href{https://lawreview.uchicago.edu/online-archive/national-pork-producers-council-v-ross-invitation-revisit-import-export-clause}{U. Chi. L. Rev. Online}).
\end{itemize}

\section{Analysis of Opinions}
\subsection{Opinion of the Court and Plurality Sections (Gorsuch, J.)}
Writing for the five-Justice majority in \textbf{Parts I, II, III, IV--A, and V}, Justice Gorsuch emphasized that Congress possesses express authority under the Commerce Clause and Supremacy Clause to enact a uniform federal pork standard and preempt state laws, but had repeatedly declined to do so (\href{https://www.law.cornell.edu/supremecourt/text/21-468}{Cornell LII}). Because petitioners disavowed any claim of purposeful discrimination, their first theory---an ``almost per se'' extraterritoriality rule---misread \textit{Baldwin}, \textit{Brown-Forman}, and \textit{Healy}, each of which invalidated protectionist price-fixing or price-affirmation statutes that forced out-of-state firms to surrender competitive advantages. In an interconnected national economy, virtually all state tax, environmental, tort, and health laws produce cross-border ripple effects; petitioners' rule would imperil centuries of traditional state police-power regulation. In \textbf{Part IV--A}, the majority explained that ``no clear line'' separates \textit{Pike} from core antidiscrimination cases: \textit{Pike} itself smoked out Arizona's protectionist cantaloupe-packing requirement, and non-protectionist invalidations have rarely extended beyond instrumentalities of interstate transportation (``Pigs are not trucks or trains'') (\href{https://www.law.cornell.edu/supremecourt/text/21-468}{Cornell LII}).

In \textbf{Part IV--B} (plurality with Thomas and Barrett, JJ.), Gorsuch argued that balancing out-of-state producers' compliance costs against California voters' moral and health interests asks judges to weigh incommensurable values---like deciding ``whether a particular line is longer than a particular rock is heavy'' (\textit{Bendix Autolite Corp. v. Midwesco Enterprises, Inc.}, 486 U.S. 888, 897 (1988) (Scalia, J., concurring in judgment)). Such policy trade-offs belong to voters, state legislatures, and Congress under the ``(wakeful) Commerce Clause.'' In \textbf{Part IV--C} (plurality with Thomas, Sotomayor, and Kagan, JJ.), Gorsuch concluded under \textit{Exxon} that petitioners failed to plead a substantial burden on interstate commerce: the Dormant Commerce Clause protects the interstate market, not particular firms or their preferred ``methods of operation,'' and compliant out-of-state producers could fill California's demand while passing costs to California consumers. In \textbf{Part IV--D} (plurality with Thomas and Barrett, JJ.), Gorsuch warned that the dissents' ``derivative harms'' approach would let a single judge invalidate state bans on horsemeat, fireworks, or plastic bags, and would violate equal state sovereignty by granting smaller States more regulatory authority than larger States (\href{https://www.law.cornell.edu/supremecourt/text/21-468}{Cornell LII}).

\subsection{Concurrence in Part (Sotomayor, J., joined by Kagan, J.)}
Justice Sotomayor joined all of Gorsuch's opinion except Parts IV--B and IV--D, writing separately to preserve \textit{Pike} balancing intact. She emphasized that courts routinely weigh disparate values and conduct means-ends tailoring across constitutional law, and agreed with Chief Justice Roberts that \textit{Pike} extends beyond discrimination and transportation cases (citing \textit{Edgar v. MITE Corp.}, 457 U.S. 624 (1982)). She voted to affirm solely on the narrower Part IV--C ground that under \textit{Exxon}, petitioners failed to clear the threshold requirement of plausibly alleging a substantial burden on interstate commerce (\href{https://www.law.cornell.edu/supremecourt/text/21-468}{Cornell LII}).

\subsection{Concurrence in Part (Barrett, J.)}
Justice Barrett joined Parts I, II, III, IV--A, IV--B, IV--D, and V, but declined to join Part IV--C. She agreed with Chief Justice Roberts that petitioners' complaint \textit{did} plausibly allege pervasive, burdensome costs felt primarily outside California. Nevertheless, she voted to affirm because \textit{Pike} balancing requires judicially comparable benefits and burdens, and California's moral interest in banning allegedly inhumane products cannot be weighed on a judicial scale against dollars and cents (\href{https://www.law.cornell.edu/supremecourt/text/21-468}{Cornell LII}).

\subsection{Concurrence in Part and Dissent in Part (Roberts, C.J., joined by Alito, Kavanaugh, and Jackson, JJ.)}
Chief Justice Roberts agreed that the Dormant Commerce Clause contains no per se extraterritoriality bar, but dissented from affirming dismissal of the \textit{Pike} claim. Stressing that a six-Justice majority preserved \textit{Pike} balancing and that a five-Justice majority found a substantial burden plausibly alleged, Roberts argued that the Ninth Circuit erred by reducing petitioners' complaint to mere ``compliance costs.'' Under \textit{Bibb v. Navajo Freight Lines, Inc.}, 359 U.S. 520 (1959), \textit{Kassel v. Consolidated Freightways Corp.}, 450 U.S. 662 (1981), and \textit{Edgar}, courts consider consequential, market-wide derivative harms---here, effectively forcing compliance on hog farmers nationwide who do not even wish to sell in California, threatening herd health, and upending established husbandry practices. Roberts would have vacated and remanded for the Ninth Circuit to conduct full \textit{Pike} balancing (\href{https://www.law.cornell.edu/supremecourt/text/21-468}{Cornell LII}).

\subsection{Concurrence in Part and Dissent in Part (Kavanaugh, J.)}
Justice Kavanaugh joined Roberts's partial dissent and wrote separately to warn that California's ``California knows best'' regulation threatened economic balkanization by leveraging its 13\% market share to export its moral and policy preferences nationwide---inviting retaliatory state embargoes over labor wages, immigration status, or abortion coverage. Notably, Kavanaugh invited future litigants to challenge Proposition 12-style production-process bans under three alternative textual provisions: the \textbf{Import-Export Clause} (Art. I, \S~10, cl. 2, questioning \textit{Woodruff v. Parham}, 8 Wall. 123 (1869)), the \textbf{Privileges and Immunities Clause} (Art. IV, \S~2, cl. 1), and the \textbf{Full Faith and Credit Clause} (Art. IV, \S~1, noting that Ohio law expressly authorized practices California banned) (\href{https://www.law.cornell.edu/supremecourt/text/21-468}{Cornell LII}).

\section{Examples: Future Applications}
\begin{enumerate}
\item \textbf{Same-Side (Upheld): State Ban on In-State Sale of Non-Compostable Product Packaging.} A large State enacts a nondiscriminatory law requiring all consumer-goods packaging sold within the State after 2032 to be recyclable or compostable (analogous to California's SB 54). Out-of-state manufacturers sue, alleging massive nationwide retooling costs. \textit{Why Same-Side:} The law applies evenhandedly to in-state and out-of-state sellers without protectionist intent, regulates only in-state sales, and---absent federal preemption by Congress---merely shifts firms' ``methods of operation'' under \textit{Exxon} and \textit{Ross} Parts III and IV--C, while serving environmental benefits that three Justices deem incommensurable under Part IV--B.

\item \textbf{Same-Side (Upheld): State Ban on In-State Sale of Fur or Foie Gras.} A State prohibits retail sales within its borders of garments made from farmed mink fur or foie gras produced by force-feeding birds, regardless of where the animals were raised. Out-of-state producers challenge the ban as an unconstitutional moral embargo. \textit{Why Same-Side:} Under \textit{Ross}, there is no ``almost per se'' extraterritoriality rule, and a State may ban in-state sales of ordinary consumer goods on nondiscriminatory moral or animal-welfare grounds, leaving dissatisfied producers to seek uniform federal legislation from Congress.

\item \textbf{Opposite-Side (Invalidated): State Law Requiring Out-of-State Retailers to Match Lowest Regional Prices.} A State enacts a statute barring out-of-state pharmaceutical distributors from selling prescription drugs in-state at prices higher than the lowest price charged in any neighboring State. \textit{Why Opposite-Side:} \textit{Ross} expressly preserved \textit{Baldwin}, \textit{Brown-Forman}, and \textit{Healy} for price-control and price-affirmation statutes that tie in-state prices to out-of-state transactions and strip out-of-state firms of competitive advantages.

\item \textbf{Opposite-Side (Invalidated): State Mandate Requiring Unique Railcar Couplers on Interstate Freight Trains.} A State passes a facially neutral safety statute requiring every freight train passing through the State to use a proprietary coupler design incompatible with the couplers mandated by neighboring States, forcing interstate railroads to decouple and reassemble trains at the state border. \textit{Why Opposite-Side:} \textit{Ross} Part IV--A (and footnote 2) explicitly distinguished ``instrumentalities of interstate transportation---trucks, trains, and the like'' (\textit{Bibb}; \textit{Southern Pacific}), where a lack of national uniformity directly impedes the physical flow of interstate goods across borders.

\item \textbf{Fence-Sitter: State Sales Ban Tied to Out-of-State Employer Labor or Reproductive-Health Policies.} A large State bars the in-state retail sale of any manufactured product made by a company that pays its out-of-state workers less than \$22 per hour---or, conversely, a large State bars the sale of goods made by companies that cover (or refuse to cover) employee abortion travel---mirroring the hypotheticals flagged in Justice Kavanaugh's partial dissent. \textit{Why Fence-Sitter:} Under \textit{Ross} Parts II--III, the law is facially nondiscriminatory and has no ``almost per se'' Dormant Commerce Clause bar. Yet unlike Proposition 12 (which regulated physical attributes of the animal producing the sold meat and asserted food-safety links), a boycott tied wholly to unrelated out-of-state corporate social policies invites renewed challenges under \textit{Pike} (where five Justices in \textit{Ross} already found pervasive out-of-state costs to be a substantial burden) as well as the Full Faith and Credit, Due Process, and Import-Export Clauses signaled by Justice Kavanaugh and \textit{Ross} Part III--B.
\end{enumerate}

\section{Critique}
\subsection{Originalist and Textualist Perspectives}
Originalists and textualists welcomed Justice Gorsuch's judicial restraint in refusing to expand the atextual Dormant Commerce Clause into a freewheeling ``roving license'' for federal judges to second-guess state legislation, praising his insistence that uniform national economic rules should come from Congress exercising its express Commerce Clause and Supremacy Clause preemption powers (\href{https://lawreview.uchicago.edu/online-archive/national-pork-producers-council-v-ross-invitation-revisit-import-export-clause}{U. Chi. L. Rev. Online}; \href{https://harvardlawreview.org/print/vol-138/federalism-rebalancing-and-the-roberts-court-a-departure-from-historical-patterns/}{Harvard Law Review, 138 Harv. L. Rev.}). At the same time, some originalist and free-market commentators criticized the fractured outcome for leaving interstate markets vulnerable to regulatory balkanization by dominant States, and highlighted Justice Kavanaugh's invitation to revisit \textit{Woodruff v. Parham} and house interstate trade protections in the Constitution's express text---specifically the Import-Export Clause---though scholars caution that stretching ``Imposts or Duties'' to cover non-tax production regulations risks creating a new atextual doctrine (\href{https://lawreview.uchicago.edu/online-archive/national-pork-producers-council-v-ross-invitation-revisit-import-export-clause}{U. Chi. L. Rev. Online}; \href{https://www.washburnlaw.edu/academics/experience/law-journal/wlj-online/volume-64/2022-03-vol-64-eating-out-of-both-troughs.html}{Washburn L.J. Online}).

\subsection{Progressive and Doctrinal Perspectives}
Progressive and animal-welfare scholars hailed \textit{Ross} as a major victory for state regulatory autonomy, democratic moral self-determination, and public health, noting that invalidating Proposition 12 would effectively force California consumers to be complicit in out-of-state farming practices they deem cruel (\href{https://www.gwlr.org/national-pork-producers-council-v-ross-a-win-for-farmed-pigs/}{GW Law Review}; \href{https://harvardlawreview.org/print/vol-137/the-dormant-commerce-clause-and-moral-complicity-in-a-national-marketplace/}{Harvard Law Review Note, 137 Harv. L. Rev.}). Doctrinal commentators, however, underscored two structural tensions: first, the 5--4 voting paradox left lower courts with deep uncertainty over how to apply \textit{Pike} after \textit{Ross}; second, as the \textit{Harvard Law Review} observed, heightened deference to state market-access regulations carries an asymmetric structural bias favoring large-economy States (such as California, New York, and Texas) capable of exporting their regulatory regimes nationwide, potentially triggering sharp interstate clashes over climate mandates, ESG rules, and social policy (\href{https://harvardlawreview.org/print/vol-137/national-pork-producers-council-v-ross/}{Harvard Law Review, 137 Harv. L. Rev. 330}).

\subsection{Underlying Competing Values}
\textit{Ross} pits \textbf{horizontal federalism and local democratic self-governance} (allowing voters in different States to decide which goods may be sold in their own markets according to local health, safety, and moral standards) against \textbf{national economic union and anti-balkanization} (protecting an integrated common market from being fragmented or effectively governed by the regulatory preferences of the largest consumer States when Congress remains silent).

\section{Key Quotations}
\begin{quote}
``Companies that choose to sell products in various States must normally comply with the laws of those various States. Assuredly, under this Court's dormant Commerce Clause decisions, no State may use its laws to discriminate purposefully against out-of-state economic interests. But the pork producers do not suggest that California's law offends this principle. Instead, they invite us to fashion two new and more aggressive constitutional restrictions on the ability of States to regulate goods sold within their borders. We decline that invitation. While the Constitution addresses many weighty issues, the type of pork chops California merchants may sell is not on that list.''
\end{quote}
\noindent --- Justice Gorsuch, Opinion of the Court, 598 U.S. at 364 (\href{https://www.law.cornell.edu/supremecourt/text/21-468}{Cornell LII}; \href{https://www.courtlistener.com/opinion/10049695/national-pork-producers-council-v-ross/}{CourtListener}).

\begin{quote}
``How is a court supposed to compare or weigh economic costs (to some) against noneconomic benefits (to others)? No neutral legal rule guides the way. The competing goods before us are insusceptible to resolution by reference to any juridical principle. Really, the task is like being asked to decide `whether a particular line is longer than a particular rock is heavy.' ''
\end{quote}
\noindent --- Justice Gorsuch, Plurality Opinion, Part IV--B, 598 U.S. at 381 (quoting \textit{Bendix Autolite Corp. v. Midwesco Enterprises, Inc.}, 486 U.S. 888, 897 (1988) (Scalia, J., concurring in judgment)) (\href{https://www.law.cornell.edu/supremecourt/text/21-468}{Cornell LII}; \href{https://www.courtlistener.com/opinion/10049695/national-pork-producers-council-v-ross/}{CourtListener}).

\begin{quote}
``In short, I vote to affirm the judgment because petitioners fail to allege a substantial burden on interstate commerce as required by \textit{Pike}, not because of any fundamental reworking of that doctrine.''
\end{quote}
\noindent --- Justice Sotomayor, Concurring in Part, 598 U.S. at 391 (\href{https://www.law.cornell.edu/supremecourt/text/21-468}{Cornell LII}; \href{https://www.courtlistener.com/opinion/10049695/national-pork-producers-council-v-ross/}{CourtListener}).

\begin{quote}
``Although \textit{Pike} is susceptible to misapplication as a freewheeling judicial weighing of benefits and burdens, it also reflects the basic concern of our Commerce Clause jurisprudence that there be `free private trade in the national marketplace.' ''
\end{quote}
\noindent --- Chief Justice Roberts, Concurring in Part and Dissenting in Part, Part I, 598 U.S. at 395 (\href{https://www.law.cornell.edu/supremecourt/text/21-468}{Cornell LII}; \href{https://www.courtlistener.com/opinion/10049695/national-pork-producers-council-v-ross/}{CourtListener}).

\begin{quote}
``The State has aggressively propounded a `California knows best' economic philosophy---where California in effect seeks to regulate pig farming and pork production in all of the United States. California's approach undermines federalism and the authority of individual States by forcing individuals and businesses in one State to conduct their farming, manufacturing, and production practices in a manner required by the laws of a different State.''
\end{quote}
\noindent --- Justice Kavanaugh, Concurring in Part and Dissenting in Part, Part I, 598 U.S. at 407 (\href{https://www.law.cornell.edu/supremecourt/text/21-468}{Cornell LII}; \href{https://www.courtlistener.com/opinion/10049695/national-pork-producers-council-v-ross/}{CourtListener}).


\end{document}
